現実のデータ収集環境では，センサーの途切れや入力漏れ，システム障害などにより欠損値が頻繁に発生する．欠損値が存在すると，統計解析や機械学習モデルの推定に歪みが生じ，分析結果の信頼性が低下するため，適切な欠損処理が不可欠である．欠損処理には削除と補完があるが，削除は情報損失を伴う．一方，平均値補完などの統計的手法は，データの非線形性や特徴量間の関係を十分に捉えられない．この課題に対し，ランダムフォレストを用いた missForest が提案されているが，すべての特徴量を一律に用いるため，寄与の小さい特徴量の影響を受けやすい．そこで本研究では，複数の特徴選択基準に基づく特徴量部分集合を用い，OOB 指標に基づいて予測を統合するアンサンブル型欠損値補完手法を提案する．
